\documentclass{article}
\usepackage[utf8]{inputenc}
\usepackage{amsmath,amsfonts,amssymb,amsthm}
\usepackage{graphicx}
\usepackage{booktabs}
\usepackage{multirow}
\usepackage{hyperref}
\usepackage[margin=1in]{geometry}
\usepackage{caption}
\usepackage{float}
\usepackage{natbib}

\renewcommand{\thetable}{S\arabic{table}}
\renewcommand{\thefigure}{S\arabic{figure}}
\renewcommand{\thesection}{S\arabic{section}}
\renewcommand{\theequation}{S\arabic{equation}}

\newtheorem{proposition}{Proposition}

\title{Supplementary Materials for\\
``\textsc{SCP}: Simulation-Based Power Calculation and Study Design for Circadian Rhythmic Analysis''}
\author{}
\date{}

\begin{document}

\maketitle


%% ====================================================================
\section{Pilot datasets and signal summaries}
\label{sec:pilot_summary}
%% ====================================================================

Table~\ref{tab:pilot_summary} summarises the pilot datasets used in
the main text figures. A complete pilot inventory (123 tissues across GTEx \citep{gtex2020atlas}, GSE98965 baboon \citep{Mure2018baboon}, GSE71620 PFC (BA11/BA47)
\citep{chen2016aging}, and GSE160521 \citep{ketchesin2021diurnal}, with
per-tissue median $\tilde r$ and rhythmic-gene counts computed on the
top-$300$ rhythmic genes) is provided as a separate supplementary
spreadsheet (\texttt{tissue\_signal\_summary.csv}). Across all entries,
$\tilde r$ is defined as the median of $r_g = A_g / \sigma_g$ over the
top-$300$ rhythmic genes ranked by cosinor F-test p-value. This matches
the gene set used to calibrate $\hat{F}_{A,\sigma}$ in the simulation
pipeline (Section~2.1 of the main text). This convention makes the
reported effect-size summaries directly comparable across pilots.

\begin{table}[!htbp]
\centering
\small
\caption{Pilot datasets used in the main-text figures. For passive
designs, collection times are empirical and sampled from the pilot
time-of-day distribution. Median effect size $\tilde r$ is the median
of $r_g = A_g/\sigma_g$ over the top-$300$ rhythmic genes from the
pilot cosinor screen at $\alpha_{\mathrm{pilot}} = 0.01$.}
\label{tab:pilot_summary}
\begin{tabular}{llccc}
\toprule
Species & Pilot dataset & $n_{\mathrm{pilot}}$ & $\tilde r$ & Used in \\
\midrule
\textit{H.\ sapiens} & GTEx Adrenal Gland & $187$ & $0.97$ & Fig.~1, Fig.~2 \\
\textit{H.\ sapiens} & GTEx Liver & $154$ & $0.48$ & Fig.~1, Fig.~4, Fig.~5 \\
\textit{H.\ sapiens} & GSE160521 Putamen control & $59$ & $1.06$ & Fig.~6 \\
\textit{H.\ sapiens} & GSE160521 Putamen schizophrenia  & $28$ & $1.02$ & Fig.~3 \\
\textit{H.\ sapiens} & GTEx Thyroid & $416$ & $0.45$ & Fig.~3 \\
\bottomrule
\end{tabular}
\end{table}


%% ====================================================================
\section{Active-design identifiability and B-invariance}
\label{sec:b_invariance}
%% ====================================================================

This section provides the algebraic derivation underlying the
$B$-invariance result stated in Section~2.1 of the main text and used in
the active-design analysis of Section~3.6. The argument extends the
cosinor noncentrality framework of \citet{zong2023circapower} to balanced
active grids and to the two-harmonic detector. Under balanced equally
spaced sampling with $B$ distinct time points and $m$ replicates per time
point ($N = Bm$), the harmonic-coefficient block of the design matrix is
diagonal and proportional to $N$ once the sampling grid satisfies the
harmonic-regression identifiability threshold.

\subsection{Single-harmonic setup}

Let $t_b = (b - 1) \cdot 24 / B$ for $b = 1, \ldots, B$ denote the
distinct sampling times of a balanced equally spaced active design over
one $24$-hour cycle, and let
$\theta_b = \omega_0 t_b = 2\pi (b - 1)/B$ be the corresponding phase.
For the single-harmonic cosinor model, the design vector is
\[
f(t)
=
\big[
1,\; \cos(\omega_0 t),\; \sin(\omega_0 t)
\big]^\top .
\]

\subsection{Orthogonality for $B \geq 3$}

The off-diagonal sums
\[
\sum_b \cos(\theta_b),
\qquad
\sum_b \sin(\theta_b),
\qquad
\sum_b \sin(\theta_b)\cos(\theta_b)
\]
vanish for $B \geq 3$ by the standard finite-sum identity for roots of
unity. At $B=2$, the sine column is identically zero, so the
single-harmonic design is rank-deficient. Thus, the active-design
identifiability threshold for the single-harmonic test is $B \geq 3$.

Using
\[
\cos^2\theta = \frac{1+\cos(2\theta)}{2},
\qquad
\sin^2\theta = \frac{1-\cos(2\theta)}{2},
\]
and the same finite trigonometric sum identity, we obtain
\[
\sum_{b=1}^{B}\cos^2\theta_b
=
\sum_{b=1}^{B}\sin^2\theta_b
=
\frac{B}{2}.
\]
With balanced replication $m_b \equiv m$ and $N=Bm$, the single-harmonic
design information matrix is
\begin{equation}
X^\top X
=
\mathrm{diag}
\big(
N,\; N/2,\; N/2
\big),
\qquad
B \geq 3.
\label{eq:supp_xtx_balanced}
\end{equation}

\subsection{Single-harmonic noncentrality and B-invariance}

Under cosinor truth,
\[
y_{g,i}
=
\mu_g
+
A_g\cos(\omega_0 t_i - \omega_0 \phi_g)
+
\varepsilon_{g,i},
\]
the harmonic-coefficient block has information $(N/2)I_2$ by
\eqref{eq:supp_xtx_balanced}. The noncentrality parameter for the
single-harmonic F-test is therefore
\begin{equation}
\lambda_g
=
\frac{N}{2}
\frac{A_g^2}{\sigma_g^2}
=
\frac{N r_g^2}{2},
\qquad
r_g = A_g/\sigma_g.
\label{eq:supp_ncp_cosinor}
\end{equation}
Thus, under balanced equally spaced active sampling, the noncentrality
depends on total sample size $N$ and effect size $r_g$, but not on the
specific allocation of samples across $B$ and $m$ once $B \geq 3$. The
same argument extends to the two-harmonic detector above its
identifiability threshold $B \geq 5$, since the harmonic columns of
\eqref{eq:supp_xtx_balanced} remain mutually orthogonal with sum of
squares $N/2$ once $B \geq 2K + 1$.


%% ====================================================================
\section{Simulation algorithm}
\label{sec:algorithm}
%% ====================================================================

Algorithm~S1 describes the data-generation procedure for a single
simulation replicate of the two-group differential setting at sample
size $N$. Expression is generated under the cosinor model
\citep{cornelissen2014cosinor}, with gene-type assignments and
endpoint-specific modifications mirroring the joint-rhythmicity strata
used by \textsc{DiffCircaPipeline} \citep{xue2023diffcircapipeline}. The
single-cohort case is recovered by setting the differential
proportions
\[
\pi_{\mathrm{DR}}
=
\pi_{\mathrm{DP}}
=
\pi_{\mathrm{DM}}
=
0
\]
and the second-group parameter modifications to no-ops.

\begin{figure}[!htbp]
\small
\hrule\vspace{4pt}
\textbf{Algorithm S1}: Two-group circadian simulation \vspace{2pt}
\hrule\vspace{4pt}
\textbf{Input}: number of genes $G$, sample size $N$, type probabilities
$\pi_R, \pi_{\mathrm{DR}}, \pi_{\mathrm{DP}}, \pi_{\mathrm{DM}}$,
empirical pilot distributions
$\hat{F}_\mu, \hat{F}_{A,\sigma}, \hat{F}_\phi$ and empirical TOD
distribution $\hat{F}_{\mathrm{TOD}}$, design type (active or passive),
phase shift $\Delta\phi$, mesor shift $\Delta\mu$, seed $s$.\\
\textbf{Output}: expression matrices
$\mathbf{Y}_1, \mathbf{Y}_2 \in \mathbb{R}^{G \times N}$ and the
ground-truth type assignment $\mathbf{T}$.
\vspace{4pt}

\begin{enumerate}
\setlength{\itemsep}{2pt}

\item \textbf{Seed} $s$. For $g=1,\ldots,G$, draw baseline parameters:
\[
\mu_g \sim \hat{F}_\mu,
\qquad
(A_g,\sigma_g) \sim \hat{F}_{A,\sigma},
\qquad
\phi_g \sim \hat{F}_\phi.
\]
The joint draw of $(A_g,\sigma_g)$ preserves the pilot's
amplitude-noise dependence.

\item \textbf{Assign gene types.} Draw
$T_g \in \{0,1,2,3,4,5\}$ from the categorical distribution
\[
P(T_g=0)
=
1-\pi_R-\pi_{\mathrm{DR}}-\pi_{\mathrm{DP}}-\pi_{\mathrm{DM}},
\]
\[
P(T_g=1)=\pi_R,
\qquad
P(T_g=2)=P(T_g=3)=\pi_{\mathrm{DR}}/2,
\qquad
P(T_g=4)=\pi_{\mathrm{DP}},
\qquad
P(T_g=5)=\pi_{\mathrm{DM}}.
\]

\item \textbf{Initialise group parameters.} Set
\[
(\mu_{g,k},A_{g,k},\sigma_{g,k},\phi_{g,k})
=
(\mu_g,A_g,\sigma_g,\phi_g)
\]
for $k\in\{1,2\}$, then modify according to $T_g$:
\[
\begin{array}{ll}
T_g=0: & A_{g,1}=A_{g,2}=0 \quad \text{(arrhythmic)},\\
T_g=1: & \text{no modification} \quad \text{(rhythmic in both)},\\
T_g=2: & A_{g,2}=0 \quad \text{(rhythmic in group 1 only)},\\
T_g=3: & A_{g,1}=0 \quad \text{(rhythmic in group 2 only)},\\
T_g=4: & \phi_{g,2}=(\phi_{g,1}+\Delta\phi)\bmod 24,\\
T_g=5: & \mu_{g,2}=\mu_{g,1}+\Delta\mu.
\end{array}
\]

\item \textbf{Reseed} $s + 7919N$. For each group $k$ and each
$i=1,\ldots,N$, draw collection times $t_{k,i}$ from
$\hat{F}_{\mathrm{TOD}}$ for passive designs, or from the equispaced
active grid otherwise. This second seeding step makes time-of-day draws
sample-size-specific while preserving the per-gene biological parameters
across $N$.

\item \textbf{Generate expression.} For each gene $g$, group $k$, and
sample $i$,
\[
y_{g,k,i}
=
\mu_{g,k}
+
A_{g,k}
\cos\!\big(
\omega_0(t_{k,i}-\phi_{g,k})
\big)
+
\varepsilon_{g,k,i},
\qquad
\varepsilon_{g,k,i}
\sim
\mathcal{N}(0,\sigma_{g,k}^2),
\]
with $\omega_0=2\pi/24$.

\item \textbf{Record ground truth.} Store
$\mathbf{T}=(T_1,\ldots,T_G)$ along with per-gene effect sizes
$r_{g,k}=A_{g,k}/\sigma_{g,k}$, phase differences, and mesor
differences.

\end{enumerate}
\vspace{2pt}\hrule
\end{figure}

The two-stage seeding in steps 1 and 4 holds the gene-type assignment
and baseline parameters fixed across sample sizes within a simulation
replicate while reseeding only the $N$-dependent quantities, including
collection times and residual noise. This isolates the effect of $N$ on
power from incidental variation in the gene-level biological
configuration.

In practice, the two-harmonic regression test is fitted by ordinary
least squares, and all reported analyses use the exact finite-sample
F-test.


%% ====================================================================
\section{Bootstrap implementation}
\label{sec:bootstrap_impl}
%% ====================================================================

The outer subject bootstrap of Section~2.3 propagates pilot sampling
variability into the projected power curve using nonparametric resampling
over pilot subjects \citep{efron1979bootstrap}. Algorithm~S2 summarises
the implementation, which is exposed as a single end-to-end entry point
in the \textsc{SCP} framework.

\begin{figure}[!htbp]
\small
\hrule\vspace{4pt}
\textbf{Algorithm S2}: Single-cohort outer-bootstrap power \vspace{2pt}
\hrule\vspace{4pt}
\textbf{Input}: pilot expression matrix
$\mathbf{Y}_{\mathrm{pilot}} \in \mathbb{R}^{G \times n_{\mathrm{pilot}}}$,
pilot collection times $\mathbf{t}_{\mathrm{pilot}}$, sample-size grid
$\{N_1,\ldots,N_J\}$, outer-resample count $B_{\mathrm{out}}$,
inner-simulation count $N_{\mathrm{sim}}^{\mathrm{in}}$, screening
threshold $\alpha_{\mathrm{pilot}}$, and FDR level $\alpha$.\\
\textbf{Output}: $(B_{\mathrm{out}} \times J)$ matrix of bootstrap
power replicates and pointwise $95\%$ confidence bands.
\vspace{4pt}

\begin{enumerate}
\setlength{\itemsep}{2pt}

\item For $b=1,\ldots,B_{\mathrm{out}}$:
\begin{enumerate}
\setlength{\itemsep}{1pt}

\item \textbf{Resample subjects.} Draw $n_{\mathrm{pilot}}$ subject
indices with replacement from $\{1,\ldots,n_{\mathrm{pilot}}\}$ to form
$\mathbf{Y}^{(b)} \in \mathbb{R}^{G \times n_{\mathrm{pilot}}}$ and the
matched time vector $\mathbf{t}^{(b)}$.

\item \textbf{Re-estimate pilot summary.} Refit the pilot summary on
$(\mathbf{Y}^{(b)},\mathbf{t}^{(b)})$ to obtain
\[
\hat{\Psi}^{(b)}
=
\big(
\hat{F}_\mu^{(b)},
\hat{F}_{A,\sigma}^{(b)},
\hat{F}_\phi^{(b)},
\hat{F}_{\mathrm{TOD}}^{(b)},
\hat{\pi}_R^{(b)}
\big).
\]

\item \textbf{Inner simulation.} For each $j$ and each
$s=1,\ldots,N_{\mathrm{sim}}^{\mathrm{in}}$, simulate a fresh
gene-by-sample matrix from $\hat{\Psi}^{(b)}$ at sample size $N_j$ via
Algorithm~S1, apply the chosen rhythmicity detector, apply BH adjustment
at $\alpha$, and record the rhythmicity discovery rate.

\item \textbf{Aggregate.} Average the discovery rate across the
$N_{\mathrm{sim}}^{\mathrm{in}}$ inner replicates to obtain
$\widehat{\mathrm{Power}}^{(b)}(N_j)$.

\end{enumerate}

\item \textbf{Summarise across outer resamples.} Compute
\[
\bar{\mathrm{Power}}(N_j)
=
B_{\mathrm{out}}^{-1}
\sum_b
\widehat{\mathrm{Power}}^{(b)}(N_j),
\]
and the pointwise $2.5$th and $97.5$th percentiles of
$\widehat{\mathrm{Power}}^{(b)}(N_j)$ across outer bootstrap replicates.

\end{enumerate}
\vspace{2pt}\hrule
\end{figure}

Default production settings are
$B_{\mathrm{out}}=50$,
$N_{\mathrm{sim}}^{\mathrm{in}}=25$, and
$\alpha_{\mathrm{pilot}}=0.01$, used for the bootstrap panels in the
main text.


%% 
%% ====================================================================
\bibliographystyle{plainnat}
\bibliography{references}

\end{document}